\clearpage\pagestyle{otMain}\setlength{\otBlocEspacement}{2.5mm}\fontsize{11}{15}\selectfont\corrigesChapitre{11}{Longueurs, masses, contenances}
\begin{blocLivre}{cor-11-1}\begin{corrige}{1}{Relations à citer}\demarche{Énoncer la relation avant le calcul.}
1 m $=100$ cm : $\res{350}$ cm. 1 kg $=1\,000$ g : $\res{750}$ g. $1\,000$ mL $=1$ L : $\res{2{,}5}$ L.
\end{corrige}\end{blocLivre}
\begin{blocLivre}{cor-11-2}\begin{corrige}{2}{Vers les millièmes}\demarche{Placer la mesure au rang des millièmes.}
1 mm $=0{,}001$ m : $\res{0{,}045}$ m. $\res{0{,}003}$ kg. $\res{0{,}007}$ L.
\end{corrige}\end{blocLivre}
\begin{blocLivre}{cor-11-3}\begin{corrige}{3}{Comparer}\demarche{Convertir dans la plus petite unité.}
$1\,200>980$ : $\res{1{,}2\ \mathrm{km}>980\ \mathrm m}$. $500<520$ : $\res{0{,}5\ \mathrm t<520\ \mathrm{kg}}$. $75>70$ : $\res{75\ \mathrm{cL}>0{,}7\ \mathrm L}$.
\end{corrige}\end{blocLivre}
\begin{blocLivre}{cor-11-4}\begin{corrige}{4}{Estimations}\demarche{Comparer aux références connues.}
Un vélo : $\res{12\ \mathrm{kg}}$. Une tasse : $\res{25\ \mathrm{cL}}$. Une salle : $\res{8\ \mathrm m}$.
\end{corrige}\end{blocLivre}
\begin{blocLivre}{cor-11-5}\begin{corrige}{5}{Additionner des mesures}\demarche{Convertir avant d’additionner.}
$1\,200+850=\res{2\,050}$ g, soit $\res{2{,}05}$ kg.
\end{corrige}\end{blocLivre}
\begin{blocLivre}{cor-11-6}\begin{corrige}{6}{Le tour du champ}\demarche{Même unité pour tous les côtés.}
$2\times(125+87{,}5)=2\times212{,}5=\res{425}$ m, soit $\res{0{,}425}$ km.
\end{corrige}\end{blocLivre}
\begin{blocLivre}{cor-11-7}\begin{corrige}{7}{La carafe et les verres}\demarche{Convertir pour diviser des entiers.}
$1{,}5$ L $=150$ cL et $150:6=\res{25}$ cL par verre.
\end{corrige}\end{blocLivre}
\begin{blocLivre}{cor-11-8}\begin{corrige}{8}{Le colis à la poste}\demarche{Tout convertir en grammes.}
Livres : $650\times3=1\,950$ g ; total $1\,950+120=\res{2\,070}$ g $=2{,}07$ kg. $2\,070>2\,000$ : $\res{\text{non}}$, il dépasse de 70 g.
\end{corrige}\end{blocLivre}
\begin{blocLivre}{cor-11-9}\begin{corrige}{9}{L’eau de la semaine}\demarche{Multiplier avant de partager.}
$2\times6=12$ bouteilles, soit $12\times1{,}5=\res{18}$ L. Par jour : $18:7\approx\res{2{,}6}$ L, car $7\times2{,}6=18{,}2$, proche de 18. La division exacte ne tombe pas juste : on arrondit.
\end{corrige}\end{blocLivre}
\begin{blocLivre}{cor-11-10}\begin{corrige}{10}{Le médicament}\demarche{Calculer la consommation quotidienne d’abord.}
Par jour : $5\times3=15$ mL. $200=15\times13+5$ : le flacon dure $\res{13}$ jours complets, et il reste 5 mL.
\end{corrige}\end{blocLivre}
\begin{blocLivre}{cor-11-11}\begin{corrige}{11}{La citerne du jardin}\demarche{Suivre les retraits dans l’ordre.}
$3$ hL $=300$ L. Arrosage : $12\times9=108$ L. Reste : $300-108-150=\res{42}$ L.
\end{corrige}\end{blocLivre}
\begin{blocLivre}{cor-11-12}\begin{corrige}{12}{Les sacs de farine}\demarche{Convertir la tonne avant de compter les sacs.}
$0{,}5$ t $=500$ kg et $25\times20=500$ : il reçoit $\res{20}$ sacs.
\end{corrige}\end{blocLivre}
\begin{blocLivre}{cor-11-13}\begin{corrige}{13}{Le ruban exact}\demarche{Travailler en centimètres.}
5 m $=500$ cm ; $500=65\times7+45$ : $\res7$ paquets et il reste $\res{45}$ cm.
\end{corrige}\end{blocLivre}
\begin{blocLivre}{cor-11-14}\begin{corrige}{14}{La tonne en question}\demarche{Comparer le total à la limite convertie.}
Sacs : $15\times35=525$ kg ; seaux : $4\times12=48$ kg ; total $\res{573}$ kg. La limite vaut 800 kg : $573<800$, c’est $\res{\text{possible}}$.
\end{corrige}\end{blocLivre}
\begin{blocLivre}{cor-11-15}\begin{corrige}{15}{L’étiquette douteuse}\demarche{Citer la relation 100 cL = 1 L.}
$75$ cL $=\res{0{,}75}$ L, et non $7{,}5$ L : l’étiquette confond les rangs. En millilitres : $\res{750}$ mL.
\end{corrige}\end{blocLivre}
